% texte.tex — f171v, vv. 765–799 (conversion initiale des travaux antérieurs)
\vers[n=765]{ΜΝ.}{Mnésiloque}{ἄγε δὴ τίς ἔσται μηχανὴ σωτηρίας;}{Voyons, quel sera l'instrument de mon salut ?}
\vers[n=766]{}{}{τίς πεῖρα, τίς ἐπίνοι᾽; ὁ μὲν γὰρ αἴτιος}{Quelle tentative? Quelle combinaison ? Car l'auteur vrai,}
\vers[n=767]{}{}{κἄμ᾽ εἰσκυλίσας εἰς τοιαυτὶ πράγματα}{celui qui m'impliqua dans cette méchante affaire,}
\vers[n=768]{}{}{οὐ φαίνετ᾽, οὔπω. φέρε τίν᾽ οὖν ἂν ἄγγελον}{n'apparaît pas encore. Allons, quel messager}
\vers[n=769]{}{}{πέμψαιμ᾽ ἐπ᾽ αὐτόν; οἶδ᾽ ἐγὼ καὶ δὴ πόρον}{lui dépêcherai-je ? (Il réfléchit)... Il me vient à l'esprit un expédient}
\vers[n=770]{}{}{ἐκ τοῦ Παλαμήδους· ὡς ἐκεῖνος, τὰς πλάτας}{tiré du \italique{Palamède}. Comme cet autre, je lancerai mes avirons}
\vers[n=771]{}{}{ῥίψω γράφων. ἀλλ᾽ οὐ πάρεισιν αἱ πλάται.}{chargés d'écritures... Mais je ne les ai pas sous la main, ces avirons...}
\vers[n=772]{}{}{πόθεν οὖν γένοιτ᾽ ἄν μοι πλάται. πόθεν; πόθεν;}{Comment donc faire pour me procurer des avirons ? Comment ?...}
\vers[n=773]{}{}{τί δ᾽ ἂν εἰ ταδὶ τἀγάλματ᾽ ἀντὶ τῶν πλατῶν}{Mais quoi, si à défaut d'avirons, j'écrivais sur ces offrandes}
\vers[n=774]{}{}{γράφων διαρρίπτοιμι; βέλτιον πολύ.}{et les jetais de tous côtés ? Cela vaut bien mieux.}
\vers[n=775]{}{}{ξύλον γέ τοι καὶ ταῦτα κἀκεῖν᾽ ἦν ξύλον.}{Elles sont de bois après tout; ceux-là aussi étaient de bois.}
\vers[n=776, indent=3em, findent=1.7999999999999998em, metre={an.}]{}{}{ὦ χεῖρες ἐμαὶ}{À l'œuvre, ô mes mains,}
\vers[n=777, indent=3em, findent=1.7999999999999998em]{}{}{ἐγχειρεῖν χρῆν ἔργῳ πορίμῳ.}{il s'agit d'un travail expédient.}
\vers[n=778, indent=3em, findent=1.7999999999999998em]{}{}{ἄγε δὴ πινάκων ξεστῶν δέλτοι,}{Allons, feuillets de mes tablettes polies,}
\vers[n=779, indent=3em, findent=1.7999999999999998em]{}{}{δέξασθε σμίλης ὁλκοὺς}{recevez les entailles du ciseau,}
\vers[n=780, indent=3em, findent=1.7999999999999998em]{}{}{κήρυκας ἐμῶν μόχθων· οἴμοι}{messagères de mes misères. Aïe,}
\vers[n=781, indent=3em, findent=1.7999999999999998em]{}{}{τουτὶ τορῶ μοχθηρόν·}{c'est pénible à percer !.. «~Ac...cours~», «~Ac...cours !~» et quel sera notre cours d'eau ?}
\vers[n=782, indent=3em, findent=1.7999999999999998em]{}{}{`χώρει', `χώρει'. ποίαν αὔλακα;}{(Il se tourne et lance les tablettes des deux côtés.)}
\vers[n=783, indent=3em, findent=1.7999999999999998em]{}{}{βάσκετ᾽ ἐπείγετε πάσας καθ᾽ ὁδοὺς}{Partez, volez par toutes les voies}
\vers[n=784, indent=3em, findent=1.7999999999999998em]{}{}{κείνᾳ ταύτᾳ· ταχέως χρή.}{avec cette nouvelle façon de faire : il faut aller vite.}
\vers[n=785, metre={an⁴}]{ΧΟ.}{Chœur}{ἡμεῖς τοίνυν ἡμᾶς αὐτὰς εὖ λέξωμεν παραβᾶσαι,}{Maintenant nous allons faire nous-mêmes notre éloge dans la parabase.}
\vers[n=786]{}{}{καίτοι πᾶς τις τὸ γυναικεῖον φῦλον κακὰ πόλλ᾽ ἀγορεύει,}{Chacun pourtant médit à l'envi de l'espèce féminine :}
\vers[n=787]{}{}{ὡς πᾶν ἐσμὲν κακὸν ἀνθρώποις κἀξ ἡμῶν ἐστιν ἅπαντα,}{nous sommes un pur fléau pour les hommes ; de nous viennent tous les maux,}
\vers[n=788]{}{}{ἔριδες νείκη στάσις ἀργαλέα λύπη πόλεμος. φέρε δή νυν,}{disputes, querelles, affreuses séditions, chagrins, guerres. Çà, voyons,}
\vers[n=789]{}{}{εἰ κακόν ἐσμεν, τί γαμεῖθ᾽ ἡμᾶς, εἴπερ ἀληθῶς κακόν ἐσμεν,}{si nous sommes un fléau, pourquoi nous épousez-vous, ce fléau que nous sommes ?}
\vers[n=790]{}{}{κἀπαγορεύετε μήτ᾽ ἐξελθεῖν μήτ᾽ ἐκκύψασαν ἁλῶναι,}{Pourquoi nous interdire de sortir, de nous laisser prendre à pousser la tête dehors ?}
\vers[n=791]{}{}{ἀλλ᾽ οὑτωσὶ πολλῇ σπουδῇ τὸ κακὸν βούλεσθε φυλάττειν;}{Pourquoi vous piquer de le garder avec un soin si jaloux, ce fléau ?}
\vers[n=792]{}{}{κἂν ἐξέλθῃ τὸ γύναιόν ποι, κᾆθ᾽ εὕρητ᾽ αὐτὸ θύρασιν,}{Arrive-t-il à votre petite femme d'être en course, et que vous la trouviez absente,}
\vers[n=793]{}{}{μανίας μαίνεσθ᾽, οὓς χρῆν σπένδειν καὶ χαίρειν, εἴπερ ἀληθῶς}{vous entrez dans des fureurs, quand il vous faudrait faire libation et vous réjouir}
\vers[n=794]{}{}{ἔνδοθεν ηὕρετε φροῦδον τὸ κακὸν καὶ μὴ κατελαμβάνετ᾽ ἔνδον.}{de ce qu'il a disparu de chez vous, le fléau, et que vous ne le trouviez pas au logis.}
\vers[n=795]{}{}{κἂν καταδάρθωμεν ἐν ἀλλοτρίων παίζουσαι καὶ κοπιῶσαι,}{Nous endormons-nous chez les autres, fatiguées de danser,}
\vers[n=796]{}{}{πᾶς τις τὸ κακὸν τοῦτο ζητεῖ περὶ τὰς κλίνας περινοστῶν.}{chacun se met après le fléau et rode autour des lits.}
\vers[n=797]{}{}{κἂν ἐκ θυρίδος παρακύπτωμεν, τὸ κακὸν ζητεῖτε θεᾶσθαι·}{Nous penchons-nous par une fenêtre, vous cherchez à le dévisager, le fléau,}
\vers[n=798]{}{}{κἂν αἰσχυνθεῖσ᾽ ἀναχωρήσῃ, πολὺ μᾶλλον πᾶς ἐπιθυμεῖ}{et si par pudeur elle se retire, chacun n'en est que plus impatient}
\vers[n=799]{}{}{αὖθις τὸ κακὸν παρακύψαν ἰδεῖν. οὕτως ἡμεῖς ἐπιδήλως}{de la voir se pencher de nouveau, ce fléau. Tant il est vrai que nous valons bien mieux}
